\eqnrefs{eq:prescribed_current_green_loss_master} 只固定平均激发数谱；要从它推出完整多量子统计，还必须说明环境的量子结构。考虑能够对角化为独立玻色自由度的线性环境，初态取这些自由度的真空，并把电子轨迹近似为不受反作用改变的给定经典电流。在线性耦合下，相互作用绘景哈密顿量可统一写成
\begin{equation}
 \boxed{
 H_I(t)
 =\sum_\lambda
 \left[f_\lambda(t)\,a_\lambda^\dagger
 +f_\lambda^*(t)\,a_\lambda\right],}
 \label{eq:prescribed-current-linear-bath-hamiltonian-dp85}
\end{equation}
其中 $\lambda$ 可以是离散模式，也可以是把频率、位置和极化连续标签合并后的 reservoir 指标；$f_\lambda(t)$ 是由给定电流与该环境自由度的场函数投影得到的复数驱动幅。定义每个自由度的累计位移幅
\begin{equation}
 \boxed{
 \alpha_\lambda
 \equiv-\frac{\ii}{\hbar}
 \int_{-\infty}^{\infty}\dd t\,f_\lambda(t).}
 \label{eq:prescribed-current-mode-displacement-dp85}
\end{equation}
若自由演化相位尚未吸收到 $f_\lambda(t)$ 中，则应把相应的 $\ee^{\ii\omega_\lambda t}$ 一并包含在该函数内。

由于\eqnrefs{eq:prescribed-current-linear-bath-hamiltonian-dp85} 在不同时刻的对易子只是复数，Magnus 展开在二阶以后终止。在这些模型条件内，任意位移幅对应的环境末态精确为多模相干态：
\begin{equation}
 \boxed{
 U_I=\ee^{\ii\Phi}\prod_\lambda D_\lambda(\alpha_\lambda),
 \qquad
 |\Psi_B^{\rm out}\rangle
 =\ee^{\ii\Phi}\bigotimes_\lambda|\alpha_\lambda\rangle.}
 \label{eq:prescribed-current-multimode-coherent-state-dp85}
\end{equation}
因此真空初态下各模式量子数彼此独立，并满足
\begin{equation*}
 N_\lambda\sim\operatorname{Poisson}(|\alpha_\lambda|^2).
\end{equation*}
连续谱极限中，\eqnrefs{eq:prescribed_current_green_loss_master} 正是这些平均占据数按频率汇总后的谱测度：
\begin{equation}
 \boxed{
 \frac{\dd P_{\rm cl}}{\dd\omega}
 =\sum_\lambda|\alpha_\lambda|^2
 \delta(\omega-\omega_\lambda),}
 \label{eq:green-loss-spectrum-mode-measure-dp85}
\end{equation}
其中对连续 reservoir 的“求和”理解为相应测度积分。

把环境获得的总能量写成
\begin{equation*}
 \Delta E_B=\sum_\lambda\hbar\omega_\lambda N_\lambda,
\end{equation*}
其特征函数 $\chi_{\Delta E_B}(s)=\langle \ee^{\ii s\Delta E_B}\rangle$ 为
\begin{equation}
 \boxed{
 \ln\chi_{\Delta E_B}(s)
 =\int_0^\infty\dd\omega\,
 \frac{\dd P_{\rm cl}}{\dd\omega}
 \left(\ee^{\ii s\hbar\omega}-1\right).}
 \label{eq:green-loss-compound-poisson-cgf-dp84}
\end{equation}
于是任意阶能量累积量都直接由 Green 损失谱给出：
\begin{equation}
 \boxed{
 \kappa_n(\Delta E_B)
 =\int_0^\infty\dd\omega\,
 (\hbar\omega)^n
 \frac{\dd P_{\rm cl}}{\dd\omega}.}
 \label{eq:green-loss-cumulants-dp84}
\end{equation}
特别地，
\begin{align}
 \langle\Delta E_B\rangle
 &=\int_0^\infty\dd\omega\,\hbar\omega
 \frac{\dd P_{\rm cl}}{\dd\omega},\\
 \operatorname{Var}(\Delta E_B)
 &=\int_0^\infty\dd\omega\,(\hbar\omega)^2
 \frac{\dd P_{\rm cl}}{\dd\omega}.
 \label{eq:green-loss-mean-variance-dp84}
\end{align}
在给定轨迹且总能量守恒的无反冲解释中，$\Delta E_B$ 等于电子能量损失。对中心纵向动量 $p_0$、群速度 $v_0$ 的入射电子，把这一经典源模型用于真实电子时还需检查由所得统计量反推的轨迹自洽性，例如显著损失窗口应满足
\begin{equation}
 \boxed{
 \epsilon_{p_\parallel,\rm pc}
 \equiv
 \frac{\sqrt{\kappa_2(\Delta E_B)+\kappa_1^2(\Delta E_B)}}{|v_0p_0|}
 \ll1,}
 \label{eq:prescribed-current-energy-selfconsistency-dp85}
\end{equation}
这里使用局域色散关系 $\delta E\simeq v_0\delta p_\parallel$ 把能量损失换成纵向动量改变量；同时还需独立要求累计横向动量转移远小于 $|p_0|$。若这些条件失效，电子反冲会改变后续耦合率，应回到量子电子--环境联合通道方程。有限温度初态、非线性环境或非 Gaussian 环境也不再给出上述真空相干态的独立 Poisson 计数。
